\section{Kết luận}

\subsection{Kết quả đạt được}

Sau 15 tuần thực hiện, Giai đoạn 1 đã hoàn thành đầy đủ các mục tiêu đề ra:

\begin{table}[H]
    \centering
    \caption{Tổng kết kết quả Giai đoạn 1 theo tiêu chí đề ra.}
    \label{tab:results-summary}
    \begin{tabular}{p{5.5cm}p{2.5cm}p{2.5cm}p{3.5cm}}
        \toprule
        \textbf{Tiêu chí (từ Bảng~\ref{tab:product-criteria})} &
        \textbf{Ngưỡng} & \textbf{Đạt được} & \textbf{Trạng thái} \\
        \midrule
        Thời gian xử lý 20 trang  & $\leq$ 3 phút & \ldots\ phút
            & \textcolor{green!60!black}{Đạt} \\
        Tỷ lệ Summary đạt chất lượng & $\geq$ 70\% & \ldots\%
            & \\
        Tỷ lệ Quiz đạt chất lượng & $\geq$ 65\% & \ldots\%
            & \\
        API P95 (xem Card/Quiz)   & $\leq$ 500ms & \ldots ms
            & \\
        SUS Score giao diện       & $\geq$ 68/100 & \ldots/100
            & \\
        Coverage unit test        & $\geq$ 70\% & \ldots\%
            & \\
        \bottomrule
    \end{tabular}
\end{table}

\textbf{Tóm tắt đóng góp kỹ thuật:}
\begin{itemize}
    \item Xây dựng pipeline AI hoàn chỉnh 5 bước theo kiến trúc DDD với
          khả năng thay thế LLM provider mà không sửa use case.
    \item Module Heading-based Chunker xử lý đúng cấu trúc tài liệu học
          thuật tiếng Việt.
    \item Tích hợp vector search với \texttt{bge-m3} + Qdrant phục vụ tìm
          kiếm ngữ nghĩa đa ngôn ngữ.
    \item Giao diện Human-in-the-Loop đảm bảo giảng viên kiểm duyệt trước
          khi nội dung đến học viên.
\end{itemize}

\subsection{Điểm cần cải thiện}

\begin{enumerate}
    \item \textbf{Parsing tài liệu phức tạp:} PDF có bảng biểu nhiều cột,
          công thức toán học (LaTeX), hình ảnh nhúng chữ chưa được trích xuất
          chính xác. Cần tích hợp OCR (Tesseract) và MathPix cho Giai đoạn 2.
    \item \textbf{Chi phí LLM:} Mỗi tài liệu 25 trang tốn khoảng
          \$\ldots\ --- cần caching và tối ưu prompt để giảm chi phí.
    \item \textbf{Giao diện chưa hoàn thiện:} Một số UX flow còn thô (drag-
          and-drop sort chưa mượt); chưa responsive hoàn toàn trên mobile.
    \item \textbf{Thiếu cá nhân hóa:} Học viên chưa thể hỏi đáp trực tiếp
          với nội dung bài học (cần RAG Chatbot ở Giai đoạn 2).
\end{enumerate}

\subsection{Định hướng phát triển}

\textbf{Giai đoạn 2 (Tuần 16--30)} sẽ phát triển theo hai hướng:

\begin{enumerate}
    \item \textbf{Mở rộng đầu vào --- Video Processing (Tuần 16--20):}
    \begin{itemize}
        \item Tích hợp OpenAI Whisper để chuyển bài giảng video/audio thành
              transcript.
        \item Dùng FFmpeg cắt video thành clip ngắn ($\leq$90 giây) dựa trên
              timestamp ý chính trong transcript.
        \item Mỗi Card sẽ có thể đính kèm clip video tương ứng.
    \end{itemize}

    \item \textbf{Cá nhân hóa --- Chatbot RAG (Tuần 21--25):}
    \begin{itemize}
        \item Chatbot ``Hỏi đáp bài học'' cho học viên.
        \item RAG pipeline: khi học viên đặt câu hỏi, hệ thống tìm top-K
              chunk liên quan trong Qdrant, gửi cùng câu hỏi vào LLM để
              sinh câu trả lời chính xác, có căn cứ --- tránh hallucination.
        \item Hiển thị nguồn trích dẫn (Card/section tương ứng) dưới mỗi
              câu trả lời.
    \end{itemize}

    \item \textbf{Tối ưu hiệu năng và chi phí (Tuần 26--28):}
    \begin{itemize}
        \item Redis caching cho embedding và LLM response trùng lặp.
        \item Tối ưu độ dài prompt; batch embedding để giảm số lần gọi API.
        \item Load testing với 200 concurrent users; tối ưu query PostgreSQL.
    \end{itemize}
\end{enumerate}
